\chapter{The CL Language}

This chapter is a condensed tour of the language. It covers what you
need in order to read and write ordinary Common Lisp, and it leaves a
good deal out. Chapter~5 lists longer books and the reference
material for when you want the rest.

Type the examples in as you go. The text after the \code{CL-USER>}
prompt is what you type, and the lines after it are what Lisp prints.

\section{Overview of CL and its Syntax}

The first thing to notice about CL is that it uses prefix
notation\index{prefix notation}. The most common thing to do, in a program or at the
REPL, is to call a function, and you do that by writing a list: the
name of the function first, then its arguments, all inside
parentheses.

\begin{verbatim}
CL-USER> (+ 2 2)
4
\end{verbatim}

This expression is the function named by the symbol \code{+},
followed by the arguments 2 and 2. Evaluating it returns 4.

\subsection{How Expressions Are Evaluated}

When CL is asked to evaluate\index{evaluation} an expression, it follows a few rules:

\begin{enumerate}
\item A number evaluates to itself:

\begin{verbatim}
CL-USER> 99
99
\end{verbatim}

\item A string, which is text in double quotes, evaluates to itself:

\begin{verbatim}
CL-USER> "Be braver -- you can't cross a chasm in two small jumps."
"Be braver -- you can't cross a chasm in two small jumps."
\end{verbatim}

\item A symbol, which is a bare name, is treated as a variable and
evaluates to the variable's value. Variables come in
Section~\ref{sec:variables}.

\item A list is treated as a call. The first element names a function
(or a macro, which we will come to). The remaining elements are the
arguments. Each argument is evaluated, left to right, and the
function is then called with the resulting values:

\begin{verbatim}
CL-USER> (expt 2 5)
32
\end{verbatim}
\end{enumerate}

Since the arguments are themselves expressions, calls nest to any
depth:

\begin{verbatim}
CL-USER> (+ (+ 2 2) (+ 3 3))
10
\end{verbatim}

\textbf{Try it:} Type these at your prompt and convince yourself that
each result makes sense:

\begin{verbatim}
(+ 2)

(+)

(user-homedir-pathname)

(get-universal-time)    ; seconds since the start of the year 1900

(search "Dr." "I am Dr. Strangelove")

(subseq "I am Dr. Strangelove"
        (search "Dr." "I am Dr. Strangelove"))
\end{verbatim}

Do not worry about getting lost among the parentheses. Your editor
indents the code to show its structure and shows you which parentheses
match, and after a short while you read Lisp by its indentation, as
you read any other language.

\subsection{Lisp Syntax Simplicity}

What you have just read is nearly the whole syntax of the language.
The one thing to add is that some operators are \emph{macros}\index{macros} or
\emph{special operators}\index{special operators} instead of functions. They are written the
same way, with the operator's name first, but they do not necessarily
evaluate all their arguments; each has its own rule. You will meet
the common ones in this chapter and their rules are easy to remember.

If you come from a language such as C, Java, JavaScript or Python, you
have been using a mixture of notations: infix for arithmetic
(\code{2 + 2}), postfix for a few things (\code{i++}), and prefix for
everything else (\code{max(a, b)}). Lisp uses the prefix form for
everything and moves the opening parenthesis to the left of the
function's name. That is the whole difference. In return there are
no precedence rules to remember, and \code{+} can take any number of
arguments.

\subsection{Turning Off Evaluation}
\label{sec:quote}

Sometimes you want an expression itself and not its value. The
special operator \code{quote}\index{quote@\texttt{quote}} returns its argument without evaluating
it:

\begin{verbatim}
CL-USER> (quote a)
A
CL-USER> (quote (+ 1 2))
(+ 1 2)
\end{verbatim}

The first returns the symbol \code{A} itself, where a bare \code{a}
would have been looked up as a variable. The second returns a list of
three items, where the unquoted form would have been a call returning
3.

This is needed so often that there is an abbreviation: a single quote
character in front of an expression.

\begin{verbatim}
CL-USER> 'a
A
CL-USER> '(+ 1 2)
(+ 1 2)
\end{verbatim}

Notice that Lisp printed the symbol in upper case. By default the
reader converts the names of symbols to upper case as it reads them,
so \code{a}, \code{A}, \code{hello} and \code{Hello} in your source
all name symbols with upper-case names. People write their code in
lower case and Lisp answers in capitals.

\subsection{Fundamental CL Data Types}

Common Lisp has the data types you know from other languages, such as
numbers, strings and arrays, and some that may be new, such as
symbols and lists. This section introduces numbers, strings, symbols
and lists. The others come later in the chapter.

CL is dynamically typed\index{dynamic typing}: values have types, and variables can
hold a value of any type. The type of a value is checked when the
program runs. You may declare types where you want the compiler to
check or optimize something, and you never have to.

\subsubsection{Numbers}

Numbers\index{numbers} in CL form a family of types: integers, ratios, floating-point
numbers and complex numbers. Most of the time you just think
``number.'' The arithmetic functions take any mixture and return a
result of the appropriate kind.

\begin{verbatim}
CL-USER> (/ 1 3)
1/3
CL-USER> (+ 1/3 2/3)
1
CL-USER> (/ 1.0 3)
0.33333334
CL-USER> (expt 2 100)
1267650600228229401496703205376
CL-USER> (sqrt -4)
#C(0.0 2.0)
\end{verbatim}

Dividing two integers gives an exact ratio\index{ratios}, so one third is really
one third, and does not get rounded until you ask for a floating-point
number. Integers have no maximum size, so arithmetic on them never
overflows.

\subsubsection{Strings}

A string\index{strings} is a sequence of characters written between double quotes.
Strings can hold any Unicode text in all current implementations.

\begin{verbatim}
CL-USER> (string-upcase "hello")
"HELLO"
CL-USER> (concatenate 'string "Hello, " "World")
"Hello, World"
\end{verbatim}

A single character is written with the prefix \code{\#\textbackslash}, as in
\code{\#\textbackslash a} for the letter a and \code{\#\textbackslash Space} for a space.

\subsubsection{Symbols}

Symbols\index{symbols} are names. You have already used some: \code{+},
\code{expt} and \code{search} are symbols that name functions.
Symbols are also data in their own right. A program can hold a
symbol in a variable, put it in a list, and compare it with another
symbol, which is very fast because two occurrences of the same name
are the very same object.

A symbol can name a function and a variable at the same time without
any conflict. \code{List}, for example, is a built-in function, and
you are free to use \code{list} as a variable name as well. There is
no need to shorten it to \code{lst}.

To get a symbol itself where Lisp would otherwise evaluate it, quote
it, as in Section~\ref{sec:quote}:

\begin{verbatim}
CL-USER> 'my-symbol
MY-SYMBOL
\end{verbatim}

Symbols whose names begin with a colon are \emph{keywords}\index{keywords}. A keyword
evaluates to itself, so it needs no quote. Keywords are used as
labels and option names throughout the language:

\begin{verbatim}
CL-USER> :red
:RED
\end{verbatim}

\subsubsection{Lists}

A list\index{lists} is written as items between parentheses. The items can be
of any type, including other lists. Since an unquoted list is a call,
you quote a list when you mean it as data:

\begin{verbatim}
CL-USER> '(a b c)
(A B C)
CL-USER> '(a (b c) d)
(A (B C) D)
\end{verbatim}

The quote covers everything inside the list.

Now look at this list:

\begin{verbatim}
(defun hello () (write-string "Hello, World!"))
\end{verbatim}

It is also a program: the function definition from Chapter~2. Take a
moment to see that it meets the description of a list. It has four
items. The first two are symbols, the third is an empty list, and the
fourth is a list of a symbol and a string.

Every CL program is made of lists like this one. That is why it is
straightforward for a CL program to read, build and transform other CL
programs, and it is the main thing that sets Lisp apart from other
languages.

Besides quoting, you can make a list by calling the function
\code{list}\index{list@\texttt{list}}, which evaluates its arguments, as every function does, and
returns a list of the results:

\begin{verbatim}
CL-USER> (list 'a 'b (+ 2 2))
(A B 4)
\end{verbatim}

\subsection{Functions}

You define a named function with \code{defun}\index{defun@\texttt{defun}}. It takes a name, a list
of parameters, and a body of one or more expressions. The value of
the last expression in the body is the value the function returns;
there is no \code{return} statement.

\begin{verbatim}
CL-USER> (defun my-first-lisp-function ()
           (list 'hello 'world))
MY-FIRST-LISP-FUNCTION
CL-USER> (my-first-lisp-function)
(HELLO WORLD)
\end{verbatim}

\code{Defun} is a macro, so it does not evaluate its arguments in the
ordinary way: the name and the parameter list are taken as written,
with no quoting.

If the first thing in the body is a string, it is kept as the
function's documentation\index{documentation strings}:

\begin{verbatim}
CL-USER> (defun square (x)
           "Return X multiplied by itself."
           (* x x))
SQUARE
CL-USER> (square 4)
16
CL-USER> (documentation 'square 'function)
"Return X multiplied by itself."
\end{verbatim}

Names in Lisp are conventionally written in lower case with hyphens
between the words.

\subsection{Global and Local Variables}
\label{sec:variables}

\subsubsection{Global Variables}

Global variables\index{global variables} are established with \code{defparameter}\index{defparameter@\texttt{defparameter}} or
\code{defvar}\index{defvar@\texttt{defvar}}:

\begin{verbatim}
CL-USER> (defparameter *todays-temp* 90)
*TODAYS-TEMP*
CL-USER> (defvar *todays-humidity* 70)
*TODAYS-HUMIDITY*
\end{verbatim}

The asterisks are part of the names. They mean nothing to Lisp. It
is a firm convention to name global variables this way so that they
stand out, and you will see why it matters in a moment.

The two differ in one way. If the variable already has a value,
\code{defvar} leaves it alone, and \code{defparameter} replaces it:

\begin{verbatim}
CL-USER> (defparameter *todays-temp* 100)
*TODAYS-TEMP*
CL-USER> (defvar *todays-humidity* 50)
*TODAYS-HUMIDITY*
CL-USER> *todays-temp*
100
CL-USER> *todays-humidity*
70
\end{verbatim}

This matters because you reload files often while you work. Use
\code{defvar} for state that should survive a reload, such as a table
your program has been filling in, and \code{defparameter} for settings
that should follow the source file.

To change the value of a variable, use \code{setf}\index{setf@\texttt{setf}}:

\begin{verbatim}
CL-USER> (setf *todays-humidity* 30)
30
CL-USER> *todays-humidity*
30
\end{verbatim}

You will also see \code{setq}\index{setq@\texttt{setq}} in older code. It does the same thing
for variables. \code{Setf} is more general, as later sections show,
so it is the one to learn.

\subsubsection{Local Variables}

Local variables\index{local variables} are introduced with \code{let}\index{let@\texttt{let}}:

\begin{verbatim}
CL-USER> (let ((a 20)
               (b 30))
           (+ a b))
50
\end{verbatim}

\code{Let} takes a list of bindings and then a body. Each binding is
a pair: a variable name and an expression for its initial value. The
body is one or more expressions, which can use the variables, and the
value of the last one is the value of the whole \code{let}. Outside
the \code{let}, the variables do not exist.

The initial values in a \code{let} are all computed before any of the
variables exists, so one binding cannot refer to another. When you
want that, use \code{let*}\index{let*@\texttt{let*}}, which makes the bindings one after
another:

\begin{verbatim}
CL-USER> (let* ((a 20)
                (b (* a 2)))
           (+ a b))
60
\end{verbatim}

Variables made by \code{let} have \emph{lexical scope}\index{lexical scope}: they can be
seen only by the code written inside the body of the \code{let}. You
can tell where a variable is visible by looking at the text of the
program.

\subsubsection{Dynamic Scope}

Global variables behave differently, and usefully so. If you bind a
global variable with \code{let}, the new value is seen not only by the
code inside the \code{let} but by every function called from there,
for as long as the \code{let} is running. Afterwards the old value is
back.

\begin{verbatim}
CL-USER> (defun report-humidity ()
           (format nil "Humidity is ~a%" *todays-humidity*))
REPORT-HUMIDITY
CL-USER> (report-humidity)
"Humidity is 30%"
CL-USER> (let ((*todays-humidity* 50))
           (report-humidity))
"Humidity is 50%"
CL-USER> (report-humidity)
"Humidity is 30%"
\end{verbatim}

This is called \emph{dynamic scope}\index{dynamic scope}, and global variables are also
known as \emph{special variables}\index{special variables} because of it. It is the clean way
to change a setting for the length of one operation. You do not have
to remember to put the old value back, even if an error interrupts the
work, and in a program with several threads the temporary value is
seen only by the thread that made it. The standard output stream, for
instance, is a special variable, so you can redirect the output of a
whole body of code by binding it.

It is also why the asterisks matter: they tell you at a glance that a
\code{let} is making a dynamic binding and not a lexical one.

\section{The List as a Data Structure}

This section presents the basic operators for working with lists.

\subsection{Accessing the Elements of a List}

The functions \code{first}\index{first@\texttt{first}} through \code{tenth} return the elements
of a list by position, and \code{nth} takes the position as a number,
counting from zero:

\begin{verbatim}
CL-USER> (first '(a b c))
A
CL-USER> (second '(a b c))
B
CL-USER> (nth 0 '(a b c))
A
CL-USER> (nth 2 '(a b c))
C
\end{verbatim}

\subsection{The ``Rest'' of the Story}

A very common pattern is to do something with the first element of a
list and then do the same with what remains. The function
\code{rest}\index{rest@\texttt{rest}} returns the list without its first element:

\begin{verbatim}
CL-USER> (rest '(a b c))
(B C)
\end{verbatim}

In older code, and in a good deal of current code, you will see
\code{first} and \code{rest} written as \code{car}\index{car@\texttt{car}} and \code{cdr}\index{cdr@\texttt{cdr}},
names that survive from the machine Lisp was first written for. They
mean the same thing.

\subsection{The Empty List}

The symbol \code{nil}\index{nil@\texttt{nil}} is the empty list, which can also be written
\code{()}. \code{Nil} evaluates to itself, so it needs no quote:

\begin{verbatim}
CL-USER> nil
NIL
CL-USER> ()
NIL
\end{verbatim}

\code{Nil} is also what Lisp uses for ``false.'' The symbol
\code{t}\index{t@\texttt{t}} is the standard value for ``true,'' but in any test,
\emph{every value other than \code{nil} counts as true}. Many
functions take advantage of this, as you will see.

The function \code{null}\index{null@\texttt{null}} tests for the empty list:

\begin{verbatim}
CL-USER> (null '(a b c))
NIL
CL-USER> (null nil)
T
\end{verbatim}

\subsection{Are You a List?}

\code{Listp} tests whether something is a list. Functions that ask a
yes-or-no question are called \emph{predicates}\index{predicates}, and by convention
their names end in \code{p}:

\begin{verbatim}
CL-USER> (listp '(pontiac cadillac chevrolet))
T
CL-USER> (listp 99)
NIL
CL-USER> (listp nil)
T
\end{verbatim}

The last result is \code{T} because \code{nil} is a list: the empty
one.

\subsection{The Conditional If}

Before going on with lists we need a way to make a choice.
\code{If}\index{if@\texttt{if}} takes three expressions: a test, a ``then'' and an ``else.''
It evaluates the test. If the result is true, it evaluates the
second expression and returns its value; otherwise it evaluates the
third and returns that. Only one of the two is ever evaluated, which
is why \code{if} cannot be a function.

\begin{verbatim}
CL-USER> (if (> 3 4)
             "yes"
             "no")
"no"
CL-USER> (if (listp '("Chicago" "Detroit" "Toronto"))
             "it is a list"
             "it is not a list")
"it is a list"
\end{verbatim}

Notice that \code{if} returns a value. In Lisp there is no division
between statements and expressions: everything returns a value and can
be used inside something else.

\subsection{Length of a List}

\code{Length}\index{length@\texttt{length}} returns the number of elements in a list:

\begin{verbatim}
CL-USER> (length '(gm ford stellantis volkswagen))
4
CL-USER> (length nil)
0
\end{verbatim}

Like most of Common Lisp, \code{length} could be written in Common
Lisp. Here is a simple version:

\begin{verbatim}
(defun our-length (list)
  (if (null list)
      0
      (+ (our-length (rest list)) 1)))
\end{verbatim}

In English: the length of an empty list is zero, and the length of any
other list is one more than the length of its rest. A function that
calls itself like this is \emph{recursive}\index{recursion}, and recursion is a natural
way to express operations on lists.

\subsection{Member of a List}

\code{Member}\index{member@\texttt{member}} tells you whether an item is in a list. If it is not,
\code{member} returns \code{nil}. If it is, \code{member} returns the
part of the list that starts with the item:

\begin{verbatim}
CL-USER> (member 'dallas '(boston san-francisco portland))
NIL
CL-USER> (member 'san-francisco '(boston san-francisco portland))
(SAN-FRANCISCO PORTLAND)
\end{verbatim}

Because anything that is not \code{nil} counts as true, the second
result works perfectly well as a ``yes'' in a test, and it carries
some useful information besides. This is a common idiom in Lisp.

\subsection{Getting Part of a List}

\code{Subseq}\index{subseq@\texttt{subseq}} returns part of a list. It takes the list, the
position to start at, and optionally the position to stop before:

\begin{verbatim}
CL-USER> (subseq '(a b c d) 1 3)
(B C)
CL-USER> (subseq '(a b c d) 1)
(B C D)
\end{verbatim}

\code{Subseq}, \code{length} and many other functions work on any
\emph{sequence}\index{sequences}, which means lists, strings and the vectors of
Section~\ref{sec:arrays}. You saw \code{subseq} applied to a string
at the start of the chapter.

\subsection{Appending Lists}

\code{Append}\index{append@\texttt{append}} takes any number of lists and returns a new list made
of all their elements:

\begin{verbatim}
CL-USER> (defparameter *slides* '(introduction welcome lists functions))
*SLIDES*
CL-USER> (append *slides* '(numbers))
(INTRODUCTION WELCOME LISTS FUNCTIONS NUMBERS)
CL-USER> *slides*
(INTRODUCTION WELCOME LISTS FUNCTIONS)
\end{verbatim}

\code{Append} did not change \code{*slides*}. Like most CL functions,
it returns a new value and leaves its arguments alone. If you want
the variable to change, say so:

\begin{verbatim}
CL-USER> (setf *slides* (append *slides* '(numbers)))
(INTRODUCTION WELCOME LISTS FUNCTIONS NUMBERS)
\end{verbatim}

\subsection{Adding Elements to a List}
\label{sec:cons}

\code{Cons}\index{cons@\texttt{cons}} adds one element to the front of a list:

\begin{verbatim}
CL-USER> (cons 'a '(b c d))
(A B C D)
\end{verbatim}

\code{Cons} is the primitive from which lists are built. A list is a
chain of pairs, called \emph{cons cells}\index{cons cells}, each holding one element
and a pointer to the next pair. \code{Cons} makes one new cell that
points at the existing list, so it is fast, and the old list is not
copied or changed. When you hear that a program ``conses\index{consing} a lot,''
it means that the program allocates a lot of memory, which the garbage
collector will later have to reclaim.

The macros \code{push}\index{push@\texttt{push}} and \code{pop}\index{pop@\texttt{pop}} combine \code{cons} and
\code{rest} with updating a variable, which lets you use a list as a
stack:

\begin{verbatim}
CL-USER> (defparameter *stack* '(b c))
*STACK*
CL-USER> (push 'a *stack*)
(A B C)
CL-USER> (pop *stack*)
A
CL-USER> *stack*
(B C)
\end{verbatim}

\subsection{Removing Elements from a List}

\code{Remove}\index{remove@\texttt{remove}} returns a new list without the occurrences of an item.
\code{Remove-if} and \code{remove-if-not} take a predicate in place of
the item:

\begin{verbatim}
CL-USER> (remove 1000 '(1 2 3 1000 4))
(1 2 3 4)
CL-USER> (remove-if-not #'evenp '(1 2 3 4 5 6))
(2 4 6)
\end{verbatim}

The notation \code{\#\q evenp} means ``the function named
\code{evenp}.'' Section~\ref{sec:functions} explains it.

\subsection{Sorting Lists}

\code{Sort}\index{sort@\texttt{sort}} sorts a sequence, given a function that says whether one
element should come before another:

\begin{verbatim}
CL-USER> (defparameter *numbers* (list 1 3 5 7 2 4 6))
*NUMBERS*
CL-USER> (setf *numbers* (sort *numbers* #'<))
(1 2 3 4 5 6 7)
CL-USER> (setf *numbers* (sort *numbers* #'>))
(7 6 5 4 3 2 1)
\end{verbatim}

\code{Sort} is an exception to the rule that functions leave their
arguments alone. For speed it is allowed to reuse the cells of the
list you give it, so afterwards the original list is in an undefined
state. Always use the value that \code{sort} returns, as above. To
keep the original, sort a copy:

\begin{verbatim}
(defun safe-sort (list predicate)
  (sort (copy-list list) predicate))
\end{verbatim}

Functions that may modify their arguments are called
\emph{destructive}\index{destructive operations}. There are only a few common ones, and their
documentation says so. Never use one on a quoted list written in your
source, because that list is a constant that belongs to the program.

\subsection{Treating a List as a Set}

\code{Union}, \code{intersection} and \code{set-difference} treat two
lists as sets. Sets have no order, so do not rely on the order of the
results:

\begin{verbatim}
CL-USER> (union '(1 2 3) '(2 3 4))
(4 1 2 3)
CL-USER> (intersection '(1 2 3) '(2 3 4))
(3 2)
CL-USER> (set-difference '(1 2 3 4) '(2 3))
(4 1)
\end{verbatim}

\subsection{Mapping a Function to a List}

When you have a list and want another list of the same length, with
each element transformed, use \code{mapcar}\index{mapcar@\texttt{mapcar}}. It takes a function and
a list, calls the function on each element, and returns a list of the
results:

\begin{verbatim}
CL-USER> (defun twice (num)
           (* num 2))
TWICE
CL-USER> (mapcar #'twice '(1 2 3 4))
(2 4 6 8)
\end{verbatim}

Given several lists, \code{mapcar} takes one argument from each:

\begin{verbatim}
CL-USER> (mapcar #'+ '(1 2 3) '(10 20 30))
(11 22 33)
\end{verbatim}

\code{Reduce}\index{reduce@\texttt{reduce}} combines all the elements of a list into one value,
using a function of two arguments:

\begin{verbatim}
CL-USER> (reduce #'+ '(1 2 3 4))
10
\end{verbatim}

\code{Mapcar}, \code{remove-if-not} and \code{reduce} are the map,
filter and reduce of other languages.

\subsection{Property Lists}

A property list\index{property lists}, or \emph{plist}, is a simple way to keep named
values. It is a list in which the odd-numbered items are names,
usually keywords, and each is followed by its value. \code{Getf}\index{getf@\texttt{getf}}
looks up a name:

\begin{verbatim}
CL-USER> (getf '(:michigan "Lansing"
                 :illinois "Springfield"
                 :pennsylvania "Harrisburg")
               :illinois)
"Springfield"
\end{verbatim}

Plists are fine for a handful of items. For large collections, use
the hash tables of Section~\ref{sec:hash-tables}.

\section{Control of Execution}

Controlling execution in CL comes down to controlling which
expressions get evaluated.

\subsection{If, When and Unless}

\code{If} takes exactly one ``then'' expression and one ``else''
expression. To do several things in either branch, group them with
\code{progn}\index{progn@\texttt{progn}}, which evaluates its arguments in order and returns the
value of the last:

\begin{verbatim}
(if (eql status :all-systems-go)
    (progn (broadcast-countdown)
           (launch-rocket))
    (progn (broadcast-apology)
           (shut-down-rocket)))
\end{verbatim}

When there is no ``else,'' use \code{when}\index{when@\texttt{when}}. It takes a test and any
number of expressions, with no need for \code{progn}.
\code{Unless}\index{unless@\texttt{unless}} is the opposite: it evaluates its body if the test is
false.

\begin{verbatim}
CL-USER> (when (> 4 3)
           (print "four is more")
           :done)

"four is more"
:DONE
CL-USER> (unless (> 4 3)
           :never)
NIL
\end{verbatim}

\subsection{Logical Operators}

\code{And}\index{and@\texttt{and}} evaluates its arguments from left to right. It stops and
returns \code{nil} as soon as one of them is false. If none is, it
returns the value of the last one. \code{Or}\index{or@\texttt{or}} stops at the first
argument that is true and returns that value:

\begin{verbatim}
CL-USER> (and (listp '(chicago detroit boston new-york))
              (listp 3.1415))
NIL
CL-USER> (or (listp '(chicago detroit boston new-york))
             (listp 3.1415))
T
CL-USER> (or nil 42 99)
42
\end{verbatim}

The last example shows a common way to supply a default:
\code{(or value default)} returns \code{value} unless it is
\code{nil}.

\code{Not}\index{not@\texttt{not}} returns \code{t} if its argument is \code{nil}, and
\code{nil} otherwise.

\subsection{Cond}

For a chain of tests, where another language would say ``if \ldots{}
else if \ldots{} else,'' use \code{cond}\index{cond@\texttt{cond}}. Each clause is a list holding
a test followed by expressions. \code{Cond} tries the tests in order,
and at the first true one it evaluates that clause's expressions and
returns. A final clause whose test is \code{t} catches everything
else:

\begin{verbatim}
CL-USER> (let ((n 4))
           (cond ((> n 10) "It's a lot")
                 ((= n 10) "It's kind of a lot")
                 (t "It's not a lot")))
"It's not a lot"
\end{verbatim}

With \code{cond} we can write our own version of \code{member}, very
close to how you would describe it in English:

\begin{verbatim}
(defun our-member (elem list)
  (cond ((null list) nil)
        ((eql elem (first list)) list)
        (t (our-member elem (rest list)))))
\end{verbatim}

``If the list is empty, the answer is no. If the item is the same as
the first element, return the list. Otherwise look in the rest.''

\subsection{Case}

To choose among fixed values, \code{case}\index{case@\texttt{case}} is more concise than
\code{cond}:

\begin{verbatim}
CL-USER> (let ((color :red))
           (case color
             (:blue "Blue is okay")
             (:red "Red is actually her favorite color")
             (otherwise "We have never heard of that!")))
"Red is actually her favorite color"
\end{verbatim}

\code{Case} compares with \code{eql}, so it works with symbols,
keywords, integers and characters, and not with strings.
Section~\ref{sec:equality} says why. \code{Ecase}\index{ecase@\texttt{ecase}} is like
\code{case} but signals an error if nothing matches, which is what
you want when the list of cases is supposed to be complete.

\subsection{Iteration}

Lisp has recursion and mapping, and it also has ordinary loops\index{iteration}.
\code{Dolist}\index{dolist@\texttt{dolist}} runs its body once for each element of a list, and
\code{dotimes}\index{dotimes@\texttt{dotimes}} runs it a given number of times, counting from zero:

\begin{verbatim}
CL-USER> (dolist (city '("Detroit" "Delft" "Tokyo"))
           (format t "Hello, ~a!~%" city))
Hello, Detroit!
Hello, Delft!
Hello, Tokyo!
NIL
CL-USER> (let ((result 0))
           (dotimes (n 10)
             (setf result (+ result n)))
           result)
45
\end{verbatim}

The \code{loop}\index{loop@\texttt{loop}} macro is a small language of its own for iteration.
It reads almost like English, and it can collect, sum, count and
filter as it goes:

\begin{verbatim}
CL-USER> (loop for n from 1 to 5 collect (* n n))
(1 4 9 16 25)
CL-USER> (loop for elem in '(4 5 6 7) sum elem)
22
CL-USER> (loop for n from 1 to 10
               when (evenp n) collect n)
(2 4 6 8 10)
\end{verbatim}

\code{Loop} can do a great deal more. The references in Chapter~5
describe it in full. To leave any of these loops early, call
\code{return}.

\section{Functions as Objects}
\label{sec:functions}

\subsection{Named Functions}

A function in CL is an object\index{functions!as objects}, like a number or a list. It can be
stored in a variable, passed to another function, and returned as a
result. To get the function object that a name refers to, write
\code{\#\q} before the name:

\begin{verbatim}
CL-USER> #'+
#<FUNCTION +>
\end{verbatim}

A function object has no printed form that could be read back in, so
Lisp prints a description between \code{\#<} and \code{>}. You will
see this notation for many kinds of object.

To call a function object, use \code{funcall}\index{funcall@\texttt{funcall}}, giving the arguments
one by one, or \code{apply}\index{apply@\texttt{apply}}, giving them as a list:

\begin{verbatim}
CL-USER> (funcall #'+ 1 2)
3
CL-USER> (apply #'+ '(1 2 3))
6
\end{verbatim}

\subsection{Functional Arguments}

You have already passed functions to \code{mapcar}, \code{sort},
\code{reduce} and \code{remove-if-not}. Functions that take other
functions as arguments are called higher-order functions\index{higher-order functions}, and CL uses
them everywhere. Your own functions qualify as soon as you define
them:

\begin{verbatim}
CL-USER> (defun add-3 (num)
           (+ num 3))
ADD-3
CL-USER> (mapcar #'add-3 '(2 3 4))
(5 6 7)
\end{verbatim}

\subsection{Anonymous Functions}

Often a function is needed in just one place, and it is not worth
giving it a name. \code{Lambda}\index{lambda@\texttt{lambda}} makes a function on the spot. It
looks like \code{defun} without the name:

\begin{verbatim}
CL-USER> (mapcar (lambda (num) (+ num 3))
                 '(2 3 4))
(5 6 7)
CL-USER> (sort (list '(4 "Buffy") '(2 "Keiko") '(1 "Judy") '(3 "Aruna"))
               (lambda (x y)
                 (< (first x) (first y))))
((1 "Judy") (2 "Keiko") (3 "Aruna") (4 "Buffy"))
\end{verbatim}

A function made by \code{lambda} can use the variables that are
visible where it is written, and it keeps them for as long as it
lives. Such a function is called a \emph{closure}\index{closures}:

\begin{verbatim}
CL-USER> (defun make-adder (n)
           (lambda (x) (+ x n)))
MAKE-ADDER
CL-USER> (defparameter *add-10* (make-adder 10))
*ADD-10*
CL-USER> (funcall *add-10* 5)
15
CL-USER> (mapcar (make-adder 100) '(1 2 3))
(101 102 103)
\end{verbatim}

Each call to \code{make-adder} returns a new function that remembers
its own \code{n}.

\subsection{Optional Arguments}

Parameters after the symbol \code{\&optional}\index{optional arguments} in a parameter list
may be left out by the caller. Each can be given a default value;
without one, the default is \code{nil}:

\begin{verbatim}
CL-USER> (defun greeting-list (&optional (username "Jake"))
           (list 'hello 'there username))
GREETING-LIST
CL-USER> (greeting-list)
(HELLO THERE "Jake")
CL-USER> (greeting-list "Joe")
(HELLO THERE "Joe")
\end{verbatim}

\subsection{Keyword Arguments}

Parameters after \code{\&key}\index{keyword arguments} are optional too, and the caller passes
them by name, in any order. They are written the same way:

\begin{verbatim}
(defun greeting-list (&key (username "Jake")
                           (greeting "how are you?"))
  (list 'hello 'there username greeting))
\end{verbatim}

\noindent
and called by putting a keyword before each value:

\begin{verbatim}
CL-USER> (greeting-list :greeting "how have you been?")
(HELLO THERE "Jake" "how have you been?")
CL-USER> (greeting-list :greeting "how have you been?" :username "Joe")
(HELLO THERE "Joe" "how have you been?")
\end{verbatim}

Keyword arguments make calls easy to read, and they let you add an
option to a function later without disturbing its existing callers.
The built-in functions use them heavily. \code{Sort}, for instance,
takes a \code{:key} function that picks out the part of each element
to compare, so the sort above could have been written
\code{(sort people \#\q< :key \#\q first)}.

A parameter after \code{\&rest}\index{rest arguments} receives all the remaining arguments
as a list, which is how you write a function that takes any number of
them:

\begin{verbatim}
CL-USER> (defun average (&rest numbers)
           (/ (reduce #'+ numbers) (length numbers)))
AVERAGE
CL-USER> (average 1 2 3 4)
5/2
\end{verbatim}

\subsection{Multiple Values}

A function can return more than one value\index{multiple values}. \code{Floor}, for
example, divides and returns both the quotient and the remainder:

\begin{verbatim}
CL-USER> (floor 7 2)
3
1
\end{verbatim}

If you use the call in the ordinary way, you get the first value and
the others are quietly dropped, so \code{(+ (floor 7 2) 10)} is 13.
To receive them all, use \code{multiple-value-bind}\index{multiple-value-bind@\texttt{multiple-value-bind}}:

\begin{verbatim}
CL-USER> (multiple-value-bind (quotient remainder)
             (floor 7 2)
           (list quotient remainder))
(3 1)
\end{verbatim}

Your own functions return several values by calling \code{values}:
\code{(values 3 1)}.

\section{Input, Output, Streams, and Strings}

Input and output in CL go through \emph{streams}\index{streams}. A stream is a
source or destination of data: the terminal, a file, a network
connection, or a string in memory. Two special variables hold the
defaults: \code{*standard-input*} and \code{*standard-output*}.

\subsection{Format}

\code{Format}\index{format@\texttt{format}} is the function you will use for most output. It takes
a destination, a control string, and arguments that are filled in
where the control string has directives. Directives start with a
tilde:

\begin{verbatim}
CL-USER> (format t "Hello there, ~a!~%" 'bob)
Hello there, BOB!
NIL
\end{verbatim}

The destination \code{t} means standard output. The directive
\code{\td a} prints the next argument in a form meant for people, and
\code{\td\%} starts a new line. \code{Format} printed its text and
then returned \code{nil}, which the REPL showed.

If the destination is \code{nil}, \code{format} prints nothing and
returns the text as a string. This is the usual way to build a string
from pieces:

\begin{verbatim}
CL-USER> (format nil "Hello there, ~a!" 'bob)
"Hello there, BOB!"
\end{verbatim}

A few more directives will cover most of what you need:
\code{\td s} prints an argument the way you would type it in a
program, \code{\td d} prints an integer, \code{\td ,2f} prints a
number with two decimal places, and \code{\td\{} \ldots{}
\code{\td\}} repeats over a list:

\begin{verbatim}
CL-USER> (format nil "~a and ~s" "a string" "a string")
"a string and \"a string\""
CL-USER> (format nil "~d items at ~,2f each" 3 4.5)
"3 items at 4.50 each"
CL-USER> (format nil "~{~a~^, ~}" '(1 2 3))
"1, 2, 3"
\end{verbatim}

There are many others. One prints numbers in English and another in
Roman numerals:

\begin{verbatim}
CL-USER> (format nil "~r" 2026)
"two thousand twenty-six"
CL-USER> (format nil "~@r" 2026)
"MMXXVI"
\end{verbatim}

\subsection{Print and Princ}

\code{Print}\index{print@\texttt{print}} outputs one object in the form that Lisp could read back
in, on a new line, and returns the object. \code{Princ}\index{princ@\texttt{princ}} outputs it in
the form meant for people. The difference shows with strings:

\begin{verbatim}
CL-USER> (print "hello")

"hello"
"hello"
CL-USER> (princ "hello")
hello
"hello"
\end{verbatim}

In each case the first thing you see is the output and the second is
the return value, printed by the REPL. \code{Print} returns its
argument, so you can wrap it around any expression in the middle of a
program to see a value go by without changing what the program does.

\subsection{Read}

\code{Read}\index{read@\texttt{read}} is the other half of the REPL: it takes text from a stream
and returns the Lisp object that the text describes. It does not
evaluate anything. The whole reader is available to your programs,
which makes Lisp's own notation a ready-made data format.

\begin{verbatim}
CL-USER> (read-from-string "(a b (c 42) \"text\")")
(A B (C 42) "text")
19
\end{verbatim}

\noindent
(The second value is the position where reading stopped.) Only read
data that you trust this way. For text from outside, read lines with
\code{read-line}\index{read-line@\texttt{read-line}}, which returns one line as a string, or use a
library for a format such as JSON, as in Chapter~4.

\subsection{File Input and Output}

To use a file you open a stream to it. The macro
\code{with-open-file}\index{with-open-file@\texttt{with-open-file}} opens the file, gives you the stream in a
variable for the length of its body, and closes the file when the body
is finished, even if it is left because of an error.

To write a file:

\begin{verbatim}
CL-USER> (with-open-file (out "/tmp/readme.txt"
                              :direction :output
                              :if-exists :supersede)
           (format out "Hello there~%")
           (format out "Second line~%"))
NIL
\end{verbatim}

\noindent
and to read it back, a line at a time:

\begin{verbatim}
CL-USER> (with-open-file (in "/tmp/readme.txt")
           (loop for line = (read-line in nil)
                 while line
                 collect line))
("Hello there" "Second line")
\end{verbatim}

Opening for input is the default, so the second form needs no options.
The \code{nil} given to \code{read-line} tells it to return
\code{nil} at the end of the file, where it would otherwise signal
an error.

\code{Probe-file}\index{probe-file@\texttt{probe-file}} tests whether a file exists.

\subsection{Pathnames}

You can name files with ordinary strings, as above. Lisp turns them
into \emph{pathname}\index{pathnames} objects, which you can take apart and combine
without cutting strings up by hand:

\begin{verbatim}
CL-USER> (defparameter *my-path* (pathname "/tmp/readme.txt"))
*MY-PATH*
CL-USER> *my-path*
#P"/tmp/readme.txt"
CL-USER> (pathname-name *my-path*)
"readme"
CL-USER> (pathname-type *my-path*)
"txt"
CL-USER> (merge-pathnames "notes.txt" *my-path*)
#P"/tmp/notes.txt"
\end{verbatim}

\section{Hash Tables, Arrays, Structures, and Classes}

Lists are convenient, and they are not the answer to everything. CL
has the other data structures you would expect.

\subsection{Hash Tables}
\label{sec:hash-tables}

A hash table\index{hash tables} maps keys to values, like a dictionary or map in other
languages. Lookup stays fast however large the table grows. Make one
with \code{make-hash-table}, look things up with \code{gethash}\index{gethash@\texttt{gethash}}, and
store them with \code{setf} of \code{gethash}:

\begin{verbatim}
CL-USER> (defparameter *ratings* (make-hash-table))
*RATINGS*
CL-USER> (setf (gethash :lisp *ratings*) :excellent)
:EXCELLENT
CL-USER> (gethash :lisp *ratings*)
:EXCELLENT
T
CL-USER> (gethash :cobol *ratings*)
NIL
NIL
\end{verbatim}

This is \code{setf}\index{setf@\texttt{setf}} doing something \code{setq} cannot. Its first
argument is not just a variable but a \emph{place}: an expression
that says where a value lives. To store into anything, you write the
expression that would fetch the value and wrap it in \code{setf}. The
same rule works for hash tables, arrays, objects and parts of lists,
so there is no separate ``put'' function to learn for each.

\code{Gethash} returns two values. The second says whether the key
was found, which lets you tell a missing key from a key whose value is
\code{nil}.

A table made with no options compares keys with \code{eql}, which
suits symbols and numbers. For string keys, ask for \code{equal}:

\begin{verbatim}
CL-USER> (defparameter *capitals* (make-hash-table :test #'equal))
*CAPITALS*
CL-USER> (setf (gethash "Michigan" *capitals*) "Lansing")
"Lansing"
CL-USER> (gethash "Michigan" *capitals*)
"Lansing"
T
\end{verbatim}

\code{Maphash} calls a function on each key and value, and
\code{hash-table-count} returns the number of entries.

\subsection{Arrays}
\label{sec:arrays}

Arrays\index{arrays} hold values in a grid of one or more dimensions, indexed by
integers from zero, with fast access to any element.
\code{Make-array} creates one and \code{aref}\index{aref@\texttt{aref}} reads an element. As
with hash tables, \code{setf} of the same expression stores one:

\begin{verbatim}
CL-USER> (defparameter *grid* (make-array '(3 3) :initial-element 0))
*GRID*
CL-USER> (setf (aref *grid* 0 0) "Dave")
"Dave"
CL-USER> (setf (aref *grid* 0 1) "John")
"John"
CL-USER> *grid*
#2A(("Dave" "John" 0) (0 0 0) (0 0 0))
\end{verbatim}

A one-dimensional array is called a \emph{vector}\index{vectors}. Vectors can be
written directly with \code{\#(} \ldots{} \code{)}, and they can be made
to grow:

\begin{verbatim}
CL-USER> (aref #(10 20 30) 1)
20
CL-USER> (defparameter *v* (make-array 0 :adjustable t :fill-pointer t))
*V*
CL-USER> (vector-push-extend :a *v*)
0
CL-USER> (vector-push-extend :b *v*)
1
CL-USER> *v*
#(:A :B)
\end{verbatim}

Strings are vectors of characters. Vectors and lists together are the
sequences, so \code{length}, \code{subseq}, \code{sort},
\code{remove}, \code{find}, \code{position} and \code{reduce} work on
all of them.

\subsection{Structures}

A structure\index{structures} is a record with named slots. One \code{defstruct}\index{defstruct@\texttt{defstruct}}
form defines the type together with a function to make instances, a
function to read each slot, and a predicate:

\begin{verbatim}
CL-USER> (defstruct point
           x
           y)
POINT
CL-USER> (defparameter *p* (make-point :x 3 :y 4))
*P*
CL-USER> (point-x *p*)
3
CL-USER> (setf (point-y *p*) 10)
10
CL-USER> *p*
#S(POINT :X 3 :Y 10)
\end{verbatim}

Structures are simple and fast. When you need more than that, use
classes.

\subsection{Classes and Methods}

The Common Lisp Object System\index{CLOS}, CLOS, is the object-oriented part of
the language. A class\index{classes} describes a kind of object: its slots, and
the classes it inherits from. \code{Defclass}\index{defclass@\texttt{defclass}} defines one:

\begin{verbatim}
(defclass box ()
  ((length :initarg :length :initform 10 :accessor box-length)
   (width  :initarg :width  :initform 10 :accessor box-width)
   (height :initarg :height :initform 10 :accessor box-height)))
\end{verbatim}

After the name comes the list of superclasses, empty here, and then
the slots. For each slot, \code{:initarg} gives the keyword used to
set it when an object is made, \code{:initform} gives its default
value, and \code{:accessor} names a function for reading the slot,
and for writing it with \code{setf}.

Behavior is defined with \emph{generic functions}\index{generic functions} and
\emph{methods}\index{methods}. A generic function is called like any other
function. What it does depends on the classes of its arguments, and
each method supplies the behavior for one case:

\begin{verbatim}
(defgeneric volume (thing)
  (:documentation "Return the volume of THING."))

(defmethod volume ((self box))
  (* (box-length self)
     (box-width self)
     (box-height self)))
\end{verbatim}

\noindent
\code{Make-instance} creates an object:

\begin{verbatim}
CL-USER> (defparameter *box* (make-instance 'box :height 2))
*BOX*
CL-USER> (volume *box*)
200
CL-USER> (setf (box-width *box*) 5)
5
CL-USER> (volume *box*)
100
\end{verbatim}

A subclass inherits slots and methods:

\begin{verbatim}
CL-USER> (defclass labeled-box (box)
           ((label :initarg :label :accessor box-label)))
#<STANDARD-CLASS COMMON-LISP-USER::LABELED-BOX>
CL-USER> (volume (make-instance 'labeled-box :label "Fragile"))
1000
\end{verbatim}

If you know Java, C++ or Python, notice what is different here.
Methods do not live inside classes. A call is written
\code{(volume box)}, like any function call, and not
\code{box.volume()}. One consequence is that you can add a method for
a class you did not write, including the built-in ones, without
touching it. Another is that a method can be chosen by the classes of
\emph{all} its arguments, not only the first:

\begin{verbatim}
CL-USER> (defgeneric add (a b))
#<STANDARD-GENERIC-FUNCTION COMMON-LISP-USER::ADD (0)>
CL-USER> (defmethod add ((a string) (b string))
           (concatenate 'string a b))
#<STANDARD-METHOD COMMON-LISP-USER::ADD (STRING STRING) {1201ACA633}>
CL-USER> (defmethod add ((a number) (b number))
           (+ a b))
#<STANDARD-METHOD COMMON-LISP-USER::ADD (NUMBER NUMBER) {1201B1BD13}>
CL-USER> (add "hello " "there")
"hello there"
CL-USER> (add 3 4)
7
\end{verbatim}

CLOS goes a good deal further than this, with multiple inheritance,
methods that run before, after and around others, and classes that can
be redefined while instances of them exist. The existing objects are
updated to match. It fits the way Lisp is developed, and it is one of
the reasons a Lisp program can stay up while it is being changed.

\section{Macros}
\label{sec:macros}

You have been using macros\index{macros} since the first \code{defun}:
\code{when}, \code{cond}, \code{dolist}, \code{loop}, \code{setf} and
\code{with-open-file} are all macros. This section shows what they
are and how to write a small one. Writing macros well is a subject
for the longer books. Reading them is something you need from the
first day.

A function receives the \emph{values} of its arguments. A macro
receives its arguments as unevaluated \emph{code}, as the lists and
symbols you typed, and returns a new piece of code to be used in their
place. It is a function from code to code, run by the compiler.
Because Lisp code is made of lists, a macro is written with the same
list operations as any other program.

Suppose \code{unless} did not exist. A function could not do its job,
because a function's arguments are all evaluated before it is called,
and the point of \code{unless} is that the body sometimes must not be.
A macro can:

\begin{verbatim}
(defmacro our-unless (test &body body)
  `(if ,test
       nil
       (progn ,@body)))
\end{verbatim}

\code{Defmacro}\index{defmacro@\texttt{defmacro}} looks like \code{defun}. The parameter \code{test}
receives the first expression in the call, and \code{\&body} collects
the rest as a list. The body builds the replacement code with the
backquote\index{backquote} notation: a backquote is like a quote, except that inside
it a comma means ``put the value of this here,'' and comma-at means
``splice the elements of this list in here.'' The result is a
template with holes.

You can ask Lisp what a macro call turns into, with
\code{macroexpand-1}\index{macroexpand-1@\texttt{macroexpand-1}}:

\begin{verbatim}
CL-USER> (macroexpand-1 '(our-unless (> 3 4) (print "no") :done))
(IF (> 3 4)
    NIL
    (PROGN (PRINT "no") :DONE))
T
\end{verbatim}

\noindent
and the new operator works like a built-in one:

\begin{verbatim}
CL-USER> (our-unless (> 3 4)
           (print "three is not more")
           :done)

"three is not more"
:DONE
\end{verbatim}

The built-in macros are no different. In SBCL:

\begin{verbatim}
CL-USER> (macroexpand-1 '(when (> 4 3) :done))
(IF (> 4 3)
    :DONE)
T
\end{verbatim}

When you meet a macro you do not know, expand it. Your editor has a
key for this (Table~\ref{tab:slime}).

This is what Chapter~1 meant by extending the language from the
inside. The built-in \code{with-open-file} is a macro that someone
wrote, and in the same way you can write your own, such as a
\code{with-connection} or a \code{define-report}, for the needs of
your program. A good rule for beginners is to
write a function whenever a function will do, and to reach for a macro
when you need to control evaluation or to remove a pattern that keeps
repeating in your code.

\section{Conditions and Error Handling}
\label{sec:conditions}

Chapter~2 showed what happens at the REPL when something goes wrong:
you land in the debugger. A finished program usually has to deal with
problems itself. In CL an error is represented by an object called a
\emph{condition}\index{conditions}, and conditions are organized into classes, such as
\code{division-by-zero}, \code{file-error} and \code{type-error}, all
of which are kinds of \code{error}.

\code{Handler-case}\index{handler-case@\texttt{handler-case}} runs an expression and says what to do if a
condition of a given class is signaled inside it. It corresponds to
try and catch in other languages:

\begin{verbatim}
CL-USER> (handler-case (/ 10 0)
           (division-by-zero ()
             :infinity))
:INFINITY
\end{verbatim}

To signal an error of your own, call \code{error}\index{error@\texttt{error}} with a message in
the style of \code{format}:

\begin{verbatim}
(defun parse-age (string)
  (let ((age (parse-integer string)))
    (when (minusp age)
      (error "An age cannot be negative: ~a" age))
    age))
\end{verbatim}

\begin{verbatim}
CL-USER> (parse-age "42")
42
CL-USER> (handler-case (parse-age "forty")
           (error (condition)
             (format t "Could not read the age: ~a~%" condition)
             nil))
Could not read the age: junk in string "forty"
NIL
\end{verbatim}

The handler names the class \code{error}, so it catches both the
error that \code{parse-integer} signaled here and the one
\code{parse-age} would signal for a negative number. The variable
\code{condition} holds the condition object, and printing it gives
its message.

To make sure something happens on the way out, whether or not there
was an error, use \code{unwind-protect}\index{unwind-protect@\texttt{unwind-protect}}. Its first expression is the
work and the rest are the cleanup:

\begin{verbatim}
CL-USER> (unwind-protect
              (progn (print "working") :result)
           (print "cleaning up"))

"working"
"cleaning up"
:RESULT
\end{verbatim}

\noindent
This is what \code{with-open-file} uses to close the file.

The condition system goes further than try and catch in one important
way. In most languages, by the time an error is caught, the code that
ran into it has been abandoned. In CL the handler runs first, while
that code is still waiting, and it can choose one of the restarts\index{restarts} that
the code has offered: use this value in place of the bad one, skip
this record, try again. The low-level code knows the possible ways to
recover and the high-level code decides which one to use. The
restarts you saw in the debugger in Section~\ref{sec:debugging} are
the same mechanism with a person making the choice. The operators to
look up when you need this are \code{handler-bind} and
\code{restart-case}.

\section{Packages}
\label{sec:packages}

A package\index{packages} is a namespace for symbols. Packages keep the names in one
library from colliding with the same names in another, and they let a
library say which of its names are meant for others to use.

There is always a current package, held in the special variable
\code{*package*}, and the REPL shows it in the prompt. A new Lisp
starts in the package \code{COMMON-LISP-USER}, or \code{CL-USER} for
short, which is a scratch area. A program of any size should have a
package of its own. \code{Defpackage}\index{defpackage@\texttt{defpackage}} defines one:

\begin{verbatim}
(defpackage :inventory
  (:use :common-lisp)
  (:export #:add-item #:*items*))
\end{verbatim}

This says that code in the \code{inventory} package can use all the
standard names of the language, which live in the package
\code{COMMON-LISP}, and that two of its own names are for the outside
world. (The \code{\#:} prefix just writes a name without creating a
symbol in the current package. You will also see keywords and strings
used here.)

Each source file then begins by saying which package it is in:

\begin{verbatim}
(in-package :inventory)

(defvar *items* nil)

(defun add-item (item)
  (push item *items*))

(defun helper ()
  :internal)
\end{verbatim}

From another package, an exported name is written with the package
name and one colon:

\begin{verbatim}
CL-USER> (inventory:add-item :widget)
(:WIDGET)
CL-USER> inventory:*items*
(:WIDGET)
\end{verbatim}

A name that was not exported can still be reached by using two
colons, as in \code{inventory::helper}. The second colon is a warning to
yourself: this name is not part of the library's interface and may
change. It is handy for poking around at the REPL and has no place in
finished code.

If you are in a different package from the one you think you are in,
Lisp will tell you that functions you have just defined do not exist.
When that happens, look at the prompt.

Package names should be distinctive, since all packages share one list
of names, and long names are tedious to type. A package can give
another package a short local nickname\index{package-local nicknames} for its own use:

\begin{verbatim}
(defpackage :report
  (:use :common-lisp)
  (:local-nicknames (:a :alexandria)))
\end{verbatim}

\noindent
after which code in \code{report} can write \code{a:flatten}. Local
nicknames came after the standard, and every implementation mentioned
in this book supports them.

\subsection{The Keyword Package}

Keywords\index{keywords}, the symbols written with a leading colon, live in a package
of their own named \code{KEYWORD}. That is why \code{:red} means the
same thing wherever it is written, whatever the current package, and
why keywords are used for option names and for labels in data.

\section{Common Stumbling Blocks}

This section goes over the things that most often trip up newcomers.

\subsection{Quotes}

A symbol without a quote is a variable. With a quote it is the symbol
itself:

\begin{verbatim}
CL-USER> (defparameter *x* 1)
*X*
CL-USER> *x*
1
CL-USER> '*x*
*X*
\end{verbatim}

A list without a quote is a call. With a quote it is the list:

\begin{verbatim}
CL-USER> (first '(+ 1 2))
+
CL-USER> (first (+ 1 2))
; Evaluation aborted on #<SIMPLE-TYPE-ERROR expected-type: LIST datum: 3>.
\end{verbatim}

In the second, \code{(+ 1 2)} was evaluated first, and \code{first}
was handed the number 3.

Numbers, strings and keywords evaluate to themselves and never need a
quote.

\subsection{Function Argument Lists}

When you define a function, the parameters go in a list of their own.
When you call it, the arguments follow the name directly:

\begin{verbatim}
CL-USER> (defun square (num)
           (* num num))
SQUARE
CL-USER> (square 4)
16
\end{verbatim}

If you are used to writing \code{square(4)}, sooner or later you will
type \code{(square (4))}. Lisp will try to call a function named
\code{4} and complain of an illegal function call.

\subsection{Symbols vs. Strings}

The symbol \code{AAA} and the string \code{"AAA"} are different kinds
of thing. A symbol is a name: there is exactly one symbol with a
given name in a given package, and comparing two symbols is a single
machine instruction. A string is a piece of text: two strings with
the same characters are still two objects.

\begin{verbatim}
CL-USER> (eql 'aaa 'aaa)
T
CL-USER> (eql "bbb" "bbb")
NIL
CL-USER> (string= "bbb" "bbb")
T
\end{verbatim}

Use symbols and keywords for things your program needs to tell apart,
such as the states of a machine or the options of a function. Use
strings for text that people read or that comes from outside. The
string operations do not apply to symbols:

\begin{verbatim}
CL-USER> (defparameter *sentence* "Jane dates only Lisp programmers")
*SENTENCE*
CL-USER> (char *sentence* 3)
#\e
CL-USER> (subseq *sentence* 0 4)
"Jane"
CL-USER> (position #\Space *sentence*)
4
CL-USER> (subseq *sentence* 11)
"only Lisp programmers"
\end{verbatim}

Remember also that the reader converts symbol names to upper case, so
the name of the symbol you type as \code{hello} is the string
\code{"HELLO"}.

\subsection{Equality}
\label{sec:equality}

CL has several functions for testing equality\index{equality}, because ``the same''
means different things for different kinds of object. A rule of
thumb:

\begin{itemize}
\item \code{eql}\index{eql@\texttt{eql}} asks whether two things are the very same object.
Use it for symbols and keywords. It is also true for two integers of
the same value and for two identical characters. It is the default
test used by \code{member}, \code{case}, \code{remove}, \code{find}
and hash tables.

\item \code{=} compares numbers by value, whatever their types.

\item \code{equal}\index{equal@\texttt{equal}} compares structure: two lists with equal elements
are \code{equal}, and so are two strings with the same characters.

\item \code{string=} compares strings, and \code{string-equal} does so
ignoring case.
\end{itemize}

\begin{verbatim}
CL-USER> (eql 3 3.0)
NIL
CL-USER> (= 3 3.0)
T
CL-USER> (eql '(a b c) (list 'a 'b 'c))
NIL
CL-USER> (equal '(a b c) (list 'a 'b 'c))
T
CL-USER> (string-equal "bbb" "BBB")
T
\end{verbatim}

The functions that use \code{eql} by default take a \code{:test}
argument to change it. Forgetting this with strings is the classic
mistake:

\begin{verbatim}
CL-USER> (member "blue" '("red" "blue" "green" "yellow"))
NIL
CL-USER> (member "blue" '("red" "blue" "green" "yellow")
                 :test #'string=)
("blue" "green" "yellow")
\end{verbatim}

Floating-point numbers that ought to be equal often differ by a tiny
amount, in Lisp as in every language. Compare them with a tolerance:

\begin{verbatim}
CL-USER> (= (+ 0.1d0 0.2d0) 0.3d0)
NIL
CL-USER> (defun nearly-equal-p (a b &optional (tolerance 1d-9))
           (< (abs (- a b)) tolerance))
NEARLY-EQUAL-P
CL-USER> (nearly-equal-p (+ 0.1d0 0.2d0) 0.3d0)
T
\end{verbatim}

\subsection{Distinguishing Macros from Functions}

Macros and functions are called with the same syntax, and it rarely
matters which is which. Assume a function unless you have a reason to
think otherwise. Most macros are easy to spot:

\begin{itemize}
\item the familiar control operators, such as \code{when},
\code{cond}, \code{case}, \code{and}, \code{or} and \code{loop};
\item names that begin with \code{def}, such as \code{defun} and
\code{defparameter};
\item names that begin with \code{with-}, such as
\code{with-open-file};
\item names that begin with \code{do}, such as \code{dolist} and
\code{dotimes}.
\end{itemize}

A small number of operators, including \code{if}, \code{let},
\code{quote}, \code{setq} and \code{progn}, are neither. They are the
special operators that the language is built on, and they behave like
macros as far as you are concerned.

If you are in doubt, ask the running Lisp:

\begin{verbatim}
CL-USER> (macro-function 'list)
NIL
CL-USER> (special-operator-p 'if)
T
\end{verbatim}

\noindent
\code{Macro-function} returns \code{nil} for anything that is not a
macro. \code{Describe} will tell you all of this at once, and your
editor's key for it is in Table~\ref{tab:slime}.

\subsection{Modifying What You Should Not}

A quoted list in your source is a constant. Lisp is allowed to store
one copy of it and share it, so changing it with a destructive
operation such as \code{sort} can change your program in ways that are
very hard to track down. If you are going to modify a list, build it
with \code{list} or \code{copy-list}, as the \code{sort} examples in
this chapter did.

Destructive operations exist for speed, and you seldom need them. A
program that builds new lists and lets the garbage collector take the
old ones is easier to get right, and usually fast enough. Write it
that way first and measure with \code{time} before you change
anything.
